% ============================================================
%  CHAPTER 5 — Discussion, Conclusions, and Future Work
% ============================================================
\chapter{Discussion, Conclusions, and Future Work}
\label{ch:conclusion}

This final chapter reflects on what was built, what was learned, and where the work
can go from here.
We begin by discussing the main design decisions and findings in the light of the
experimental results, then situate the work within the broader literature, and close
with conclusions and a road-map for future development.

% ============================================================
\section{Discussion}
\label{sec:discussion}
% ============================================================

\subsection{VLM-Driven Task Decomposition in Practice}

The central architectural bet of this project was that a zero-shot cloud VLM could
replace hand-crafted finite state machines as the decision engine for a physical robot.
The experimental results support this bet.
In successful trials Gemini 2.0 Flash consistently selected the correct skill sequence
without any robot-specific fine-tuning: it understood that you cannot hand over an
object you have not yet grasped, that you should search before grasping an invisible
object, and that the task is finished once the gripper opens toward the user.

The structured JSON output format --- enforced by the system prompt with the constraint
\emph{``Respond ONLY with valid JSON, no markdown, no prose''} --- was essential.
Early iterations without this constraint produced free-text replies that required
fragile regex parsing and failed roughly 20\% of the time.
The current design reduces parse failures to negligible levels, with the retry logic
handling any edge cases.

Crucially, the affordance feedback loop proved to be what transformed the VLM from a
brittle oracle into a robust planner.
Without it, a single VLM error --- say, selecting \texttt{grab\_bottle} when the gripper
is already closed --- would leave the agent stuck in an infinite retry loop.
With it, the failure message is injected into the next prompt, and the model corrects
itself, typically within one additional iteration.
This closed-loop self-correction is one of the most practically valuable design patterns
to emerge from the project.

\subsection{3-D Localisation: Where Noise Matters Most}

The RANSAC cylinder fitting pipeline achieved sub-5\,mm positional stability on stationary
objects, which turned out to be more than accurate enough for the gripper to close around
a 500\,ml bottle.
The 5-frame averaging reduced the dominant noise source (single-frame depth jitter from
the Xtion sensor's structured-light pattern) to an acceptable level.

The main residual failure mode was \emph{insufficient depth coverage}: thin objects or
bottles at the far edge of the detection bounding box sometimes produced fewer than 15
valid depth points, causing the RANSAC fitting to be rejected and the grasp attempt to
abort.
This suggests that the confidence threshold for RANSAC acceptance (currently a minimum
inlier count) should be adaptive rather than fixed.

\subsection{The Torso Descent: Unconventional but Effective}

The torso-descent vertical approach --- using the prismatic torso joint for the final
approach instead of a Cartesian arm path --- was arguably the most unusual engineering
decision in the project, and the one that made the difference between a system that
worked occasionally and one that worked reliably.

The key insight is geometric: the TIAGo torso is a pure world-Z-aligned prismatic joint.
Lowering it while the arm is extended above the object guarantees a straight vertical
descent with no IK computation, no singularity risk, and no lateral deviation.
MoveIt is used only for the higher-level joint-space move to position the gripper above
the object; the final centimetre-precision descent is purely kinematic.

The main cost is speed: the torso descends at 3\,cm/s for safety, making the approach
phase 3--4 seconds regardless of the required travel distance.
For a demonstration scenario this is acceptable; for a productive robot in a time-critical
environment it would not be.
The planned v6 implementation (see Section~\ref{sec:future_work}) will address this with
a full MoveJ/MoveL trajectory generator.

\subsection{The Split-Process Architecture: Constraint as Opportunity}

The split-process architecture --- YOLO running as an HTTP microservice on the host, the
agent running in the Docker container --- was initially an unwanted consequence of the
Python 3.6 constraint.
In retrospect it turned out to be a clean design.

The separation of concerns is genuine: the inference layer (YOLO model, GPU if available,
Ultralytics API) is completely decoupled from the ROS control layer.
Upgrading the detector requires only restarting the microservice; the agent code is
untouched.
The HTTP interface is also trivially testable: any developer can POST a JPEG to
\texttt{localhost:5001/detect} and inspect the JSON response without launching ROS.
The 18\,ms round-trip latency (including HTTP overhead) is negligible compared to the
agent's 1--4\,s VLM query time.

\subsection{Limitations and Honest Caveats}

The system has several limitations that are important to state clearly.

\paragraph{No object permanence.}
Objects are tracked only within the current camera frame.
If the bottle moves out of view during a grasp attempt, the system has no memory of its
last known position and must restart with a \texttt{search\_with\_head} call.
A scene graph updated across loop iterations would address this at modest implementation
cost.

\paragraph{Fixed grasp orientation.}
The grasp orientation quaternion is fixed for downward vertical approach.
This works for upright cylindrical objects on flat surfaces but fails for objects lying
on their side or oriented at an angle.
Extending the pipeline to estimate the object's principal axis from the RANSAC cylinder
fit would generalise grasping to a much wider range of objects.

\paragraph{Internet dependency.}
All reasoning requires a live connection to the Gemini API.
Network interruptions degrade the agent to a single fallback skill (\texttt{go\_to\_safe\_pose}).
A locally deployed VLM would eliminate this dependency at the cost of requiring a
GPU-equipped companion machine.

\paragraph{COCO-80 class restriction.}
YOLO11n detects only the 80 COCO classes.
Requesting ``bring me the blue mug'' would fail if the mug is not close enough to
``cup'' in class space.
This restriction disappears when GPU resources permit the use of open-vocabulary detectors
such as Grounding DINO~\cite{liu2023gdino}.

\subsection{Comparison with Related Work}

Table~\ref{tab:related_comparison} situates this work relative to the most prominent
VLA systems in the literature.

\begin{table}[H]
\centering
\caption{Comparison with related vision-language-action robot systems.}
\label{tab:related_comparison}
\begin{tabularx}{\linewidth}{lXXX}
\toprule
\textbf{System} & \textbf{VLM / LLM} & \textbf{Skills} & \textbf{Grasp approach} \\
\midrule
SayCan \cite{ahn2022saycan}   & PaLM + value fn  & 551 skills & Learned policy \\
RT-2 \cite{brohan2023rt2}     & PaLM-E (fine-tuned) & End-to-end & Learned trajectory \\
VILA \cite{lin2024vila}        & VILA VLM (fine-tuned) & End-to-end & Learned trajectory \\
This work                      & Gemini 2.0 Flash (zero-shot) & 9 skills & RANSAC + MoveIt \\
\bottomrule
\end{tabularx}
\end{table}

The critical difference is training data: SayCan collected 68,000 real-robot episodes;
RT-2 fine-tuned a 55-billion-parameter model on a joint vision-language-action dataset;
VILA required multi-stage training on curated embodied datasets.
This work requires \emph{no robot-specific training data at all} --- only calibration
of the grasp offsets, which takes roughly 30 minutes on a new platform.

The trade-off is that the nine hand-crafted skills are less general than learned
policies: the system cannot smoothly interpolate to novel objects or motion profiles.
But for the narrow task domain targeted here (single-object pick-and-place + handover in
a known environment), the zero-shot approach achieves comparable task success rates at
a fraction of the development cost.

\subsection{Lessons Learned}

Five lessons stand out from the development process.

\begin{enumerate}

  \item \textbf{Iterative development on real hardware is non-negotiable.}
        The grasp pipeline required five versions to reach reliable performance.
        Each version was motivated by a specific failure mode identified on the physical
        robot that no simulation would have caught: sensor noise patterns, gripper
        mechanical play, MoveIt singularity regions near table height.

  \item \textbf{Sensor noise must be actively managed.}
        Single-frame RANSAC estimates are too noisy for sub-centimetre grasping.
        Five-frame averaging reduced positioning errors from $\sim$15\,mm to $<$5\,mm
        at negligible latency cost.

  \item \textbf{Structured prompts are critical for reliable VLM integration.}
        The difference between a reliable agent and an unreliable one was largely in the
        system prompt: strict JSON schema, explicit behavioural rules, and affordance
        failure injection.
        Prompt engineering is as important as code engineering in VLM-based systems.

  \item \textbf{Safety should be designed in from the start.}
        The affordance gating system prevented several potentially damaging arm movements
        during early development.
        Bolting on safety after the fact is much harder and less reliable.

  \item \textbf{Architectural constraints can drive good design.}
        The Python 3.6 constraint led to the HTTP microservice architecture, which turned
        out to be a clean separation of concerns that is beneficial regardless of the
        constraint.
        Constraints force creativity.

\end{enumerate}

% ============================================================
\section{Conclusions}
\label{sec:conclusions}
% ============================================================

This thesis presented the design, implementation, and evaluation of a complete
Vision--Language--Action agent for the PAL Robotics \tiago{} mobile manipulator.
The system accepts natural-language commands via voice or text, reasons about the scene
using a cloud VLM, and executes physical manipulation skills --- all without any
robot-specific training data.

The six main contributions are:

\begin{enumerate}

  \item \textbf{A complete, deployable VLA agent} integrating Google Gemini 2.0 Flash,
        YOLO11n object detection, and nine physical skills into a coherent
        perceive--reason--act loop that runs within the constraints of a Python~3.6
        / ROS~Melodic Docker environment.

  \item \textbf{A robust 3-D object localisation pipeline} using RANSAC cylinder fitting
        on RGB-D point clouds with 5-frame temporal averaging, achieving sub-5\,mm
        positional stability for stationary objects.

  \item \textbf{A split-process YOLO microservice} that overcomes the Python~3.6
        deployment constraint by running YOLO11n on the host machine and communicating
        via HTTP, achieving $\approx$18\,ms inference with a clean architectural separation
        between the AI inference layer and the ROS control layer.

  \item \textbf{An iterative grasping pipeline} that evolved through five versions, each
        addressing a specific real-robot failure mode, culminating in the torso-descent
        vertical approach that achieves reliable single-object grasps.

  \item \textbf{A handover skill} that detects a person using YOLO and depth sensing,
        plans a collision-free arm trajectory, and presents a held object at an ergonomic
        height --- completing the full ``find--grasp--deliver'' task chain demonstrated in
        the lab.

  \item \textbf{A safety-first architecture} with per-skill affordance gating, consecutive
        failure counting, and automatic safe-mode recovery to prevent robot damage during
        failures.

\end{enumerate}

The results demonstrate that a zero-shot VLM-based agent can achieve reliable multi-step
manipulation on a real mobile manipulator without robot-specific fine-tuning.
The modular architecture --- where perception, reasoning, and execution are cleanly
separated --- makes it straightforward to upgrade individual components independently.
This work establishes a solid foundation for more capable embodied AI agents as better
models, faster inference, and richer sensor suites become available.

% ============================================================
\section{Future Work}
\label{sec:future_work}
% ============================================================

Several natural extensions would meaningfully increase the system's capabilities.

\subsection{Industrial-Style Trajectory Generation (v6)}

The torso-descent approach is effective but slow ($\sim$3--4\,s approach phase) and
mechanically unusual.
Version 6 of the grasping pipeline will replace it with a proper MoveJ/MoveL trajectory
generator using Linear Segment with Parabolic Blends (LSPB) velocity profiles, providing
smooth, time-optimal motions that generalise to any object type and produce more natural
robot behaviour.

\subsection{Navigation Integration}

Navigation is implemented (Section~\ref{sec:architecture}) but not yet wired into the
agent's skill library.
Closing this gap requires:
\begin{itemize}[noitemsep]
  \item registering \texttt{NavigateToSkill} in the skill registry;
  \item integrating a named-location manager (save / load / query positions);
  \item adding multi-room tasks: ``find the bottle on the kitchen shelf''.
\end{itemize}

\subsection{Object Permanence and Semantic Memory}

A scene graph updated across loop iterations would allow the agent to remember where
objects were last seen and reason about persistence.
This requires modest engineering: a Python dictionary mapping object labels to their
last known 3-D positions in the base frame, updated at each successful detection.

\subsection{Local VLM Deployment}

Replacing the Gemini API with a locally deployed model (LLaVA~1.6, VILA, or InternVL)
would eliminate the internet dependency and reduce latency.
A GPU-equipped workstation connected via ROS would enable inference at $<$200\,ms per
query.

\subsection{Adaptive Grasp Orientation}

Extending the grasping pipeline to estimate the object's principal axis from the RANSAC
cylinder fit and adapt the approach orientation accordingly would enable grasping of
non-cylindrical objects lying in arbitrary orientations --- significantly expanding the
range of manipulation tasks the system can handle.

\subsection{Open-Vocabulary Detection}

When GPU resources are available, replacing YOLO11n with an open-vocabulary detector
such as Grounding DINO~\cite{liu2023gdino} or YOLOE~\cite{wang2025yoloe} would remove
the restriction to COCO-80 classes and enable fully open-ended task specifications:
\emph{``bring me the red mug''} rather than \emph{``bring me the bottle''}.

% ============================================================
\section{Final Remarks}
% ============================================================

The demo that motivated this work --- a user saying \emph{``Hello TIAGo, I am thirsty,
can you do something about that?''} and the robot searching, grasping a bottle, finding
the user, and handing it over --- is a small but concrete step toward robots that are
genuinely useful as everyday companions.

What makes it noteworthy is not the individual components, each of which has been
demonstrated in isolation many times before, but the \emph{integration}: a physical
robot, in a real room, reasoning about what to do with a cloud brain, and executing it
with classical robotic control.
The gap between language understanding and physical action is closed --- not by
end-to-end learning, but by a carefully engineered bridge between the two worlds.

As VLMs become faster, cheaper, and more capable, the reasoning quality will improve
for free.
The architecture built here is designed to benefit directly from those improvements
without requiring a redesign.
